\oldpage[IV-2]{145}

\subsection{Propiedades de transitividad.}

\begin{proposition}[6.6.1]
\label{IV.2.s6.6.1-fr}
Para un preesquema localmente noetheriano $Z$, designemos por $\boldsymbol{P}(Z)$ una
cualquiera de las propiedades siguientes:
\begin{enumerate}
\item[\textup{a)}] $Z$ es un preesquema de Cohen-Macaulay.
\item[\textup{b)}] $Z$ verifica $(S_n)$.
\item[\textup{c)}] $Z$ es regular.
\item[\textup{d)}] $Z$ verifica $(R_n)$.
\item[\textup{e)}] $Z$ es normal.
\item[\textup{f)}] $Z$ es reducido.
\end{enumerate}
Sean entonces $X$, $Y$, $Z$ tres preesquemas localmente noetherianos,
$f:X\to Y$, $g:Y\to Z$ dos morfismos.
\begin{enumerate}
\item[\textup{(i)}] Supongamos que $f$ y $g$ son planos y que para todo
$y\in f(X)$ \textup{(resp.} todo $z\in g(Y)$), el preesquema $f^{-1}(y)$
\textup{(resp.} $g^{-1}(z)$) verifica $\boldsymbol{P}$. Entonces $h=g\circ f$ es plano y para todo
$z\in h(X)$, $h^{-1}(z)$ verifica $\boldsymbol{P}$.
\item[\textup{(ii)}] Supongamos verificadas las condiciones siguientes: $f$ es
fielmente plano, $h=g\circ f$ es plano, para todo $y\in Y$, el
preesquema $f^{-1}(y)$ verifica $\boldsymbol{P}$, y para todo $z\in h(X)$ el
preesquema $h^{-1}(z)$ verifica $\boldsymbol{P}$. Entonces $g$ es plano, y para todo
$z\in g(Y)$, $g^{-1}(z)$ verifica $\boldsymbol{P}$.
\end{enumerate}
\end{proposition}

\begin{enumerate}
\item[(i)] Ya se sabe (2.1.6) que $h$ es plano; por otra parte, para todo
      $z\in Z$,
$f_z=f\otimes 1_{\mathbf{k}(z)}:h^{-1}(z)\longrightarrow g^{-1}(z)$ es plano (2.1.4), y para todo $y\in g^{-1}(z)$, $f_z^{-1}(y)$ es
isomorfo a $f^{-1}(y)$ por transitividad de las fibras (\textbf I, 3.6.4). La
conclusión resulta entonces, respectivamente, de (6.3.5, (ii)), (6.4.2),
(6.5.2, (ii)), (6.5.3, (ii)), (6.5.4, (ii)) y (3.3.5, (ii)).
\item[(ii)] Ya se sabe que $g$ es plano (2.2.13); además, para todo
      $z\in Z$, $f_z$ es fielmente plano (2.2.13). La conclusión
      resulta entonces, respectivamente, de (6.3.5, (i)), (6.4.2),
(6.5.2, (i)), (6.5.3, (i)), (6.5.4, (ii)) y (2.1.13).
\end{enumerate}

\begin{remark}[6.6.2]
\label{IV.2.s6.6.2-fr}
Supongamos que $h$ es plano y $f$ fielmente plano, y supongamos que para
todo $z\in Z$, $h^{-1}(z)$ es \emph{de coprofundidad} $\leqslant n$ (5.7.1);
entonces resulta del razonamiento de (6.6.1, (ii)) y de (6.3.2) que
$g^{-1}(z)$ es \emph{de coprofundidad} $\leqslant n$ para todo $z\in g(Y)$ y
que para todo $y\in Y$, $f^{-1}(y)$ es \emph{de coprofundidad} $\leqslant n$.
\end{remark}
